% IEEE conference paper: Few-Shot Law Acquisition: What Transfers Across Physics Laws, Measured
% Compile: pdflatex ieee_paper.tex   (figures in ../figs)
\documentclass[conference,10pt]{IEEEtran}

\usepackage{graphicx}
\usepackage{amsmath,amssymb}
\usepackage{booktabs}
\usepackage[table]{xcolor}
\usepackage{url}
\usepackage[margin=1in]{geometry}

\begin{document}

\title{Few-Shot Law Acquisition: What Transfers Across Physics Laws, Measured}

\author{\IEEEauthorblockN{Sehaj Randhir Singh}
\IEEEauthorblockA{\textit{Department of Electrical and Computer Engineering,}\\
\textit{NYU Tandon School of Engineering (independent researcher)}\\
Brooklyn, NY, USA\\
sehajrsinghs@gmail.com}}

\maketitle

\begin{abstract}
Few-shot transfer in physics ML usually means: same equation, more parameters. We transfer across laws: a generalist that has seen five laws adapts to a held-out sixth with a quarter of the data at 2.9$\times$ lower error than a from-scratch specialist. Then we decompose why, with ablations that overturned the intended narrative: the quantity vocabulary is the dominant carrier (removing it is 10$\times$ worse), and the physics residual -- helpful in general -- constrains trajectory specialization at tiny budgets (removing it improves trajectory error 7$\times$ while answers stay flat). Both findings are per-seed, committed, and re-runnable.
\end{abstract}

\begin{IEEEkeywords}
few-shot learning; transfer learning; physics-informed machine learning; transformers
\end{IEEEkeywords}

\section{Introduction}
What does a model that has learned several physics laws know that a from-scratch learner does not? This paper measures the transfer and then decomposes its mechanism with ablations, reporting the decomposition as measured -- including the direction that surprised the authors.

\section{Protocol}
Five training laws, one held-out law, budgets from 24 to 96 samples. Conditions: generalist with real equation signature, dummy-signature control, from-scratch specialist. Ablations: no vocabulary stream, no physics residual, frozen body.

\section{Results}
\begin{table}[!t]
\caption{Few-shot law acquisition (25\% budget) and the ablation
decomposition. The vocabulary is the dominant carrier; the physics
residual constrains trajectory specialization at tiny budgets.}
\label{tab:fewshot}
\centering
\begin{tabular}{lc}
\toprule
Condition & median rel-MAE \\
\midrule
Generalist (real signature) & 0.1971046675358058 \\
Dummy control & 0.23896641936098253 \\
From-scratch specialist & 0.5803608226154111 \\
Ablation: no vocabulary & 2.021912348943419 \\
Ablation: no physics & 0.21398293623972417 \\
Ablation: frozen body & 0.19808596781154184 \\
\bottomrule
\end{tabular}
\end{table}

At the 25\% budget the generalist is 2.9$\times$ better than the specialist. Removing the vocabulary is catastrophic (10$\times$); removing the physics residual frees trajectory specialization 7$\times$ while answers stay flat -- the residual constrains curves at tiny budgets before data can specialize them.

\begin{thebibliography}{1}
\bibitem{brown2020language} T.~Brown et al., ``Language models are few-shot learners,'' NeurIPS 2020.
\end{thebibliography}

\end{document}
